% Part: history
% Chapter: biographies 
% Section: alan-turing 

\documentclass[../../../include/open-logic-section]{subfiles}

\begin{document}

\olfileid{his}{bio}{tur} 

\olsection{آلان تورنغ}

وُلد آلان تورنغ في مايدا فيل بلندن في 23 يونيو 1912. ويُعد أبا لعلم
الحاسوب النظري. بدأ اهتمام تورنغ بالعلوم الفيزيائية والرياضيات في سن
مبكرة. غير أن اهتماماته في صباه لم تلق تمثيلًا حسنًا في مدارسه، حيث
كان التركيز على الأدب والدراسات الكلاسيكية. لذلك كان أداؤه المدرسي
ضعيفًا ووبخه كثير من معلميه.

\olphoto{turing-alan}{Alan Turing}

التحق تورنغ بكلية كينغز في كامبريدج طالبًا جامعيًا، ودرس فيها
الرياضيات. وفي عام 1936 طوّر ما يسمى اليوم آلة تورنغ، في محاولة لتعريف
مفهوم الدالة القابلة للحساب تعريفًا دقيقًا وإثبات عدم قابلية مسألة
القرار للحسم. لكن ألونزو تشرش سبقه إلى النتيجة، إذ أثبتها بواسطة حساب
لامبدا الخاص به. ومع ذلك نُشرت ورقة تورنغ مع الإشارة إلى نتيجة تشرش.
ودعا تشرش تورنغ إلى برينستون، فأقام فيها بين عامي 1936 و1938 ونال
الدكتوراه بإشرافه.

على الرغم من اهتمام تورنغ بالمنطق، ظلت اهتماماته المبكرة بالعلوم
الفيزيائية بارزة. ووُظفت مهاراته العملية أثناء خدمته في قسم تحليل
الشفرات البريطاني في بلتشلي بارك خلال الحرب العالمية الثانية. وكان
تورنغ شخصية محورية في كسر الشفرة التي استخدمتها اتصالات البحرية
الألمانية، وهي شفرة إنيغما. وقد أتاحت خبرته في الإحصاء والتعمية، إلى
جانب إدخال الآلات الإلكترونية، للفريق كسر الشفرة بإنشاء آلة لفكها
سميت «بومب». وأسهمت أفكاره أيضًا في إنشاء أول حاسوب إلكتروني قابل
للبرمجة في العالم، كولوسوس، الذي استُخدم هو الآخر في بلتشلي بارك لكسر
شفرة لورنتس الألمانية.

كان تورنغ مثليًا. ومع ذلك، تقدم في عام 1942 للزواج من جوان كلارك، إحدى
زميلاته في بلتشلي بارك، ثم فسخ الخطبة لاحقًا وأخبرها بأنه مثلي. وكانت
له علاقات عاطفية عدة طوال حياته، مع أن الأفعال الجنسية بين الرجال
كانت آنذاك جرائم يعاقب عليها القانون في المملكة المتحدة. وفي عام 1952
سطا صديق لعشيقه آنذاك على منزله، وحين قدم تورنغ بلاغًا إلى الشرطة أقر
بوجود علاقة مثلية، ظنًا منه أن الحكومة كانت في سبيلها إلى إضفاء
الشرعية على الأفعال المثلية. ولم يكن ذلك صحيحًا، فوُجهت إليه تهمة
الفحش الجسيم. واختار تورنغ، بدل دخول السجن، علاجًا هرمونيًا يخفض
الرغبة الجنسية. وعُثر عليه ميتًا في 8 يونيو 1954 من جرعة زائدة من
السيانيد، وكان ذلك على الأرجح انتحارًا. ومنحته الملكة إليزابيث الثانية
عفوًا ملكيًا عام 2013.

\begin{reading}
للاطلاع على سيرة شاملة لآلان تورنغ، انظر~\citet{Hodges2014}. وقد ألهمت
حياته وأعماله مسرحية~\emph{كسر الشفرة}، التي أُنتجت للتلفزيون عام 1996
وقام ديريك جاكوبي فيها بدور تورنغ. كما أن~\emph{لعبة التقليد}، وهو
فيلم مرشح لجائزة الأوسكار من بطولة بنديكت كامبرباتش وكيرا نايتلي،
مستوحى بصورة غير صارمة من حياة آلان تورنغ ومرحلة عمله في بلتشلي بارك
\citep{Imitation2014}.

تقدم~\citet{Radiolab2012} عدة حلقات صوتية عن حياة تورنغ وعمله. كما
يمكن مشاهدة وثائقي بي بي سي هورايزن~\emph{الحياة الغريبة للدكتور
  تورنغ وموته} على الإنترنت~\citep{Sykes1992}. ويقدم~\citep{Theelen2012}
مقطعًا قصيرًا لآلة تورنغ عاملة مصنوعة من مكعبات ليغو، أُنجزت تكريمًا
لمئوية ميلاد تورنغ عام 2012.

ورقة تورنغ الأصلية عن آلات تورنغ ومسألة القرار هي
\citet{Turing1937}.
\end{reading}

\end{document}
